\section{审稿级复核与可复现性设计}

数学建模竞赛论文的可信度不仅来自模型形式，还来自“数据从哪里来、结果如何生成、结论能否被第三方复算”三个问题。本节把论文结果拆成可验证的证据链，并说明在预测更新、人员缺勤和效率变化时如何复用同一套框架。与附录中的逐行明细相比，本节关注的是复核方法本身和指标之间的逻辑关系，因此属于正文分析而不是结果堆叠。

\subsection{结果证据链与哈希式管理}

设原始数据文件为$D$，预处理结果为$D'$，模型参数为$\Theta$，求解输出为$Y$，图表和表格为$G$，最终PDF为$P$。本文的生成关系可以表示为
\begin{equation}
 D\xrightarrow{\mathrm{clean}}D'\xrightarrow{\Theta}Y
 \xrightarrow{\mathrm{render}}G\xrightarrow{\mathrm{compile}}P.
\end{equation}
每条正文结论$C_k$都应能回溯到$Y$中的字段，每幅图$G_j$都应记录数据来源和绘图脚本，每个PDF版本都应记录编译入口及页数。因而，复核不应只检查“文件存在”，还应检查路径属于项目目录、输入文件未被意外替换，以及输出时间戳和结果摘要相互一致。

\begin{table}[htbp]
\centering
\caption{论文结论的证据链结构}
\label{tab:evidence-chain}
\begin{tabular}{p{.22\textwidth}p{.26\textwidth}p{.36\textwidth}}
\toprule
结论类别 & 直接证据 & 独立复核方式\\
\midrule
日内最少人数 & 每日班次与人数、目标值 & 重新计算小时容量和库存守恒\\
服务水平增量 & 截止前处理量、早间到货量 & 检查$P_d^{(0,15)}-A_d^{(0,12)}\ge0$\\
月度最少编制 & 581人下界与出勤矩阵 & 检查峰值覆盖、23个工日和滚动窗口\\
敏感性结论 & 三个效率情景的结果摘要 & 重新运行参数扰动并比较目标值\\
图表结论 & 图形数据源和绘图脚本 & 检查图中总量与结果JSON一致\\
\bottomrule
\end{tabular}
\end{table}

在工程实现中，结果清单采用“声明路径—计算摘要—校验状态”的三元组。若某项结果只有截图而没有结构化数据，或只有结构化数据而没有生成脚本，则只能作为展示材料，不能作为高强度结论的唯一证据。本文将完整逐小时表格保留在生成目录和Excel支撑文件中，正文使用统计图、边际量和典型日表格表达规律，从而兼顾可读性与可追溯性。

\subsection{典型日剖面与极端日剖面}

全月平均图适合发现总体形状，但容易掩盖第12日等峰值日的约束压力。为此，本文同时采用典型日、峰值日和低谷日三类剖面。对任意日期$d$，定义标准化到货曲线
\begin{equation}
 \hat a_{dh}=\frac{a_{dh}}{A_d},
\end{equation}
它消除了日总量差异，可以单独比较到货时序；再定义累计到货曲线
\begin{equation}
 F_d(h)=\frac{\sum_{\tau=0}^{h}a_{d\tau}}{A_d},
\end{equation}
用于观察货物在日内的集中程度。若$F_d(12)$较大，则早间服务截止更可能成为主导约束；若晚间仍有较大增量，则起点为16的班次和日末清零约束更重要。

\begin{figure}[htbp]
\centering
\includegraphics[width=.86\textwidth]{figures/arrival_heatmap.pdf}
\caption{全月到货时序：按日期和小时识别早晚高峰}
\label{fig:arrival-body}
\end{figure}

图\ref{fig:arrival-body}不直接给出排班答案，而是说明为什么必须采用逐小时流量模型。相同的日总量如果沿小时轴平移，理论总量下界不变，但可行班次组合和16时服务压力会发生变化。由此可见，日总量是规模变量，时序形状是结构变量，两者应在数据分析阶段分别保留。

\begin{table}[htbp]
\centering
\caption{三类日期的诊断维度}
\label{tab:day-profile}
\begin{tabular}{p{.18\textwidth}p{.25\textwidth}p{.42\textwidth}}
\toprule
日期类型 & 主要特征 & 优先检查指标\\
\midrule
峰值日 & 日总量和峰值人数同时较高 & 班次覆盖、最大库存、16时截止余量\\
早间压力日 & $F_d(12)$较高但总量未必最高 & 早班人数、低效率小时错峰\\
低谷日 & 总量较低且存在较多空闲时段 & 班次数量约束、冗余能力和人员轮休\\
晚间压力日 & 20--23时仍有较大到货 & 起点16班次、日末库存清零\\
\bottomrule
\end{tabular}
\end{table}

\subsection{约束活跃性与瓶颈定位}

对一个已经求得的可行解，可以把约束分为活跃和非活跃两类。容量约束的松弛量为$e_{dh}$，截止约束的松弛量为$S_d$，覆盖约束的松弛量为
\begin{equation}
 H_d=\sum_iw_{id}-r_d.
\end{equation}
当$e_{dh}=0$时，该小时的作业能力是瓶颈；当$S_d$接近0时，16时服务承诺缺少缓冲；当$H_d$接近0时，月度出勤在该日没有可替代人员。将三种松弛量放在同一张诊断表中，可以区分日内瓶颈和月度瓶颈。

\begin{figure}[htbp]
\centering
\includegraphics[width=.86\textwidth]{figures/backlog_profile.pdf}
\caption{库存轨迹：用上升段和下降段定位货物流瓶颈}
\label{fig:backlog-body}
\end{figure}

图\ref{fig:backlog-body}中库存上升并不意味着模型失败，只要库存随后被消化且在23时归零，说明排班允许合理的暂存。真正的风险是库存上升后没有足够的下降区间，或在16时服务截止附近仍有较大早间货物未处理。这个判别也解释了为什么本文同时报告库存曲线和截止增量，而不是只报告总人日。

\subsection{月度方案的公平性与替代性}

问题三的基本目标是最小化招工人数，但同一个581人规模可能存在多种出勤矩阵。为了比较这些可行方案，可增加二级评价指标。令$n_i^{\mathrm{night}}$为工人$i$承担晚班的次数，$\bar n^{\mathrm{night}}$为其平均值，则夜班公平性可以用
\begin{equation}
 J_{\mathrm{fair}}=\sum_{i=1}^{N}\left(n_i^{\mathrm{night}}-\bar n^{\mathrm{night}}\right)^2
\end{equation}
衡量；出勤切换次数可以用
\begin{equation}
 J_{\mathrm{switch}}=\sum_{i=1}^{N}\sum_{d=1}^{29}|w_{i,d+1}-w_{id}|
\end{equation}
衡量。实际求解可先固定$N=581$和每日覆盖约束，再最小化$J_{\mathrm{fair}}+\eta J_{\mathrm{switch}}$，从而在不增加编制的前提下改善员工体验。

\begin{figure}[htbp]
\centering
\includegraphics[width=.86\textwidth]{figures/roster_heatmap.pdf}
\caption{人员--日期矩阵：检查休息日分散性与覆盖结构}
\label{fig:roster-heat-body}
\end{figure}

图\ref{fig:roster-heat-body}的价值在于把抽象的滚动窗口约束变成可视结构：若休息日形成大面积垂直空洞，说明排班可能在某些日期过度依赖少数人员；若休息日呈分散条带，则更符合月度轮休的管理直觉。图形不能替代逐人检查，但能快速发现矩阵生成或编号分配中的结构性错误。

\subsection{预测更新与模型稳定性}

在真实业务中，逐小时到货量会随预测更新而变化。设第$t$次预测为$a^{(t)}_{dh}$，第$t+1$次预测为$a^{(t+1)}_{dh}$，则重排触发条件可以写为
\begin{equation}
 \max_{d,h}\frac{|a^{(t+1)}_{dh}-a^{(t)}_{dh}|}{\max(1,a^{(t)}_{dh})}>\varepsilon,
\end{equation}
或直接比较早间需求和峰值需求是否越过预设阈值。触发后，先重新求解受影响日期的问题一、问题二，再把新的需求向量传给问题三；若固定人员仍可覆盖，则只调整日期--班次分配，尽量不改变已公布的休息日。

\begin{table}[htbp]
\centering
\caption{滚动重排的四级响应策略}
\label{tab:rolling}
\begin{tabular}{p{.18\textwidth}p{.24\textwidth}p{.43\textwidth}}
\toprule
级别 & 触发信号 & 响应方式\\
\midrule
一级 & 小幅预测修订 & 只重算处理量与库存，不改变人员\\
二级 & 早间需求明显增加 & 重算当日班次起点和低效率小时\\
三级 & 覆盖需求超过581人 & 启用临时工或调整固定编制\\
四级 & 连续多日效率下降 & 重新估计能力参数并评估设备/培训投入\\
\bottomrule
\end{tabular}
\end{table}

该策略体现了模型的稳定性原则：数据变化只影响必要的决策层，硬约束和证据链结构保持不变。每次重排都应生成新的结果摘要、图表和编译记录，并在发布门禁中检查旧PDF没有被误当作新结果。这样，论文框架同时适用于竞赛论文的最终提交和实际运营的滚动计划。

\subsection{独立复算清单}

最终提交前，复核人员可按以下顺序执行：先检查原始数据的720条联合键，再检查30个日模型的5个非空班次；随后逐小时回算处理量和库存，验证日末为零；接着核对问题二的截止约束与低效率工时总数；最后检查问题三的581名工人、23个工日、连续8日窗口和每日覆盖。任何一项失败都应阻止PDF发布，而不是只在结果表中标注“待修正”。

\begin{table}[htbp]
\centering
\caption{终稿发布前的最小复核集合}
\label{tab:release-check}
\begin{tabular}{p{.26\textwidth}p{.28\textwidth}p{.30\textwidth}}
\toprule
复核对象 & 必须满足的条件 & 失败后的处理\\
\midrule
数据 & 日期、小时、单位和联合键完整 & 回到预处理阶段\\
日内流量 & 库存非负且日末清零 & 重新求解对应日期\\
服务截止 & 16时前处理量不小于早间到货 & 调整早班或效率参数\\
月度人员 & 581人、23工日、窗口和覆盖均通过 & 重新求解人员矩阵\\
图表与PDF & 来源路径、页数和摘要一致 & 阻止发布并重新编译\\
\bottomrule
\end{tabular}
\end{table}

\subsection{本节小结}

通过证据链、约束活跃性、典型日剖面和滚动重排机制，本文把“得到一个可行解”提升为“得到一个可以解释、复算和持续更新的方案”。正文保留决定性图表和方法论，逐行原始明细则由程序和Excel生成物承载；二者通过结果字段、文件路径和独立校验关联起来。该组织方式既减轻正文阅读负担，又确保评审者能够从任意一个核心数字追溯到原始数据和求解过程。

\clearpage
\subsection{结果摘要与决策接口}

为了让论文结论能够直接服务于排班决策，本文将最终输出压缩为“规模、时序、服务、编制”四类接口。规模接口回答需要多少人，时序接口回答这些人应在何时到岗，服务接口回答截止承诺是否有余量，编制接口回答月度固定员工能否覆盖每日需求。四类接口分别对应模型中的目标值、班次起点、库存与截止松弛量、出勤矩阵，彼此之间通过同一份结果文件连接。

\begin{table}[htbp]
\centering
\caption{终稿面向决策者的四类结果接口}
\label{tab:decision-summary}
\begin{tabular}{p{.19\textwidth}p{.24\textwidth}p{.40\textwidth}}
\toprule
接口 & 核心输出 & 决策用途\\
\midrule
规模接口 & 10791、11951人日；峰值537、581人 & 制定日计划和效率改善目标\\
时序接口 & 5个班次起点及各班次人数 & 安排早班、晚班和机动岗位\\
服务接口 & 库存曲线、截止余量、小时余量 & 识别16时服务风险与瓶颈小时\\
编制接口 & 581名员工、23个工日及覆盖矩阵 & 发布月度轮休表并处理临时缺勤\\
\bottomrule
\end{tabular}
\end{table}

在论文评审中，这种接口化表达有助于区分“模型结论”和“原始明细”：前者应在正文中解释因果关系和适用边界，后者应在附件或生成物中支持复算。对于本题，最重要的结论不是某一张长表，而是问题一到问题三的逻辑传递，以及效率下降5\%时峰值需求上升到611人的风险信号。后续若更换数据集，只需重新生成四类接口，正文的模型结构和复核规范仍可复用。

\clearpage
\subsection{评审者阅读路径与结论定位}

本文建议按“问题定义—约束建模—数值结果—独立复核—管理含义”的顺序阅读。问题重述和数据分析明确了为什么必须保留小时维度；问题一至问题三给出逐层增加现实约束后的数学模型；库存、利用率和敏感性图表解释最优值的来源；最后的证据链与发布清单回答结果是否能够复算。这样的阅读路径可以避免把问题三的581人误解为问题一的日均人数，也避免用附件长表替代对模型机制的理解。

\begin{table}[htbp]
\centering
\caption{正文中的问题定位索引}
\label{tab:reader-path}
\begin{tabular}{p{.22\textwidth}p{.23\textwidth}p{.37\textwidth}}
\toprule
审阅问题 & 对应正文位置 & 应关注的证据\\
\midrule
为什么不能只除以总工时？ & 数据分析、问题一 & 到货热力图、库存守恒、连续班次约束\\
为什么问题二人数更高？ & 问题二、综合决策实验 & 低效率小时、16时截止、增量图\\
581人是否为最优编制？ & 问题三 & 峰值下界、23工日、连续8日窗口\\
结果是否真的可执行？ & 检验与实现、复核章节 & 独立回算、覆盖矩阵、发布门禁\\
换一批数据能否复用？ & 滚动重排、证据链 & 参数接口、结果注册、生成脚本\\
\bottomrule
\end{tabular}
\end{table}

从竞赛评审角度看，模型的说服力来自三种一致性。第一是维度一致性：货物以件计、人数以人计、目标以人日计，所有换算都在能力约束中显式出现；第二是逻辑一致性：每一个结果都能由前一层输出解释，问题三不重新假设一个独立需求；第三是证据一致性：正文数字、图表、结构化结果和最终PDF使用同一轮生成记录。三种一致性同时成立时，论文才具备从公式到结论的闭环，而不仅是格式完整的排版文件。
